\subsection{Eigenvalue and Singular-Value Interlacing}

\begin{theorem} \label{thm:rank-one-psd-eigenvalue-interlacing} \hlt{Rank-One Positive Perturbation Interlacing}\\
Let $A,E\in\operatorname{Mat}(n,n,\mathbb{C})$ be Hermitian, with $E\succeq0$ and $\operatorname{rank}(E)\leq1$.\\
Then $\lambda_i(A)\leq\lambda_i(A+E)\leq\lambda_{i+1}(A)\quad(1\leq i<n)$, $\lambda_n(A)\leq\lambda_n(A+E)$\\
Equivalently, $\lambda_i(A-E)\leq\lambda_i(A)\leq\lambda_{i+1}(A-E)\quad(1\leq i<n)$, $\lambda_n(A-E)\leq\lambda_n(A)$.
\begin{proof}
Let $\mathbf{u}_j$ and $\mathbf{v}_j$ be orthonormal eigenvectors of $A$ and $A+E$, respectively, in increasing eigenvalue order.\\
For the lower bound, set $\mathcal{U}=\operatorname{span}(\mathbf{u}_i,\ldots,\mathbf{u}_n)$, $\mathcal{V}=\operatorname{span}(\mathbf{v}_1,\ldots,\mathbf{v}_i)$. Since $\dim\mathcal{U}+\dim\mathcal{V}=n+1$, choose a unit $\mathbf{x}\in\mathcal{U}\cap\mathcal{V}$. Then $\lambda_i(A)\leq \mathbf{x}^*A\mathbf{x}\leq \mathbf{x}^*(A+E)\mathbf{x}\leq\lambda_i(A+E)$ by Theorem \ref{thm:rayleigh-quotient-spectral-interval} and $E\succeq0$.\\
For the upper bound, set $\mathcal{U}=\operatorname{span}(\mathbf{u}_1,\ldots,\mathbf{u}_{i+1})$, $\mathcal{V}=\operatorname{span}(\mathbf{v}_i,\ldots,\mathbf{v}_n)$, $\mathcal{W}=\ker(E)$. Here $\dim(\mathcal{U}\cap\mathcal{V})\geq2$ and $\dim\mathcal{W}\geq n-1$, so choose a unit $\mathbf{x}\in\mathcal{U}\cap\mathcal{V}\cap\mathcal{W}$. Thus $\lambda_i(A+E)\leq \mathbf{x}^*(A+E)\mathbf{x} =\mathbf{x}^*A\mathbf{x} \leq \lambda_{i+1}(A)$.\\
The statements for $A-E$ follow by applying the proved result to $A-E$ and $E$.\\
\hlt{Structural Commentary}: Interlacing comes from test vectors in overlapping spectral subspaces.\\
\hlt{Key Insight}: Rank one leaves a large kernel, so the upper-bound vector can avoid the perturbation.
\end{proof}
\end{theorem}

\begin{corollary} \label{cor:rank-two-one-positive-one-negative-interlacing} \hlt{Rank-Two Indefinite Interlacing}\\
Let $A,H\in\operatorname{Mat}(n,n,\mathbb{C})$ be Hermitian, $n\geq2$. Suppose $\operatorname{rank}(H)\leq2$ and $H$ has at most one positive and one negative eigenvalue. Then
\begin{align}
\lambda_1(A+H)\leq\lambda_2(A),\qquad
\lambda_{i-1}(A)\leq\lambda_i(A+H)\leq\lambda_{i+1}(A)\quad(2\leq i<n),
\qquad
\lambda_{n-1}(A)\leq\lambda_n(A+H).
\nonumber
\end{align}
\begin{proof}
Write $H=E_+-E_-$ with $E_\pm\succeq0$ and $\operatorname{rank}(E_\pm)\leq1$. Set $B=A-E_-$; then $A=B+E_-$ and $A+H=B+E_+$. For $2\leq i<n$, rank-one interlacing gives $\lambda_{i-1}(A)\leq\lambda_i(B)\leq\lambda_i(A+H)\leq\lambda_{i+1}(B)\leq\lambda_{i+1}(A)$.\\
The endpoints follow similarly: $\lambda_1(A+H)\leq\lambda_2(B)\leq\lambda_2(A)$, and $\lambda_{n-1}(A)\leq\lambda_n(B)\leq\lambda_n(A+H)$.\\
\hlt{Structural Commentary}: Such an $H$ is a difference of two rank-one PSD perturbations.\\
\hlt{Key Insight}: Interlace once with $E_-$, then once with $E_+$.
\end{proof}
\end{corollary}

\begin{theorem} \label{thm:cauchy-principal-submatrix-interlacing} \hlt{Cauchy Interlacing: One Principal Deletion}\\
Let $A=
\begin{bmatrix}
B & \mathbf{y}\\
\mathbf{y}^* & c
\end{bmatrix}
\in\operatorname{Mat}(n+1,n+1,\mathbb{C})$ be Hermitian, with $B\in\operatorname{Mat}(n,n,\mathbb{C})$.\\
Then, for $1\leq i\leq n$, $\lambda_i(A)\leq\lambda_i(B)\leq\lambda_{i+1}(A)$.
\begin{proof}
Let $\mathbf{u}_j$ and $\mathbf{v}_j$ be orthonormal eigenvectors of $A$ and $B$, respectively, in increasing eigenvalue order. Identify $\mathbb{C}^n$ with $\{(\mathbf{z},0)\}\subseteq\mathbb{C}^{n+1}$, so $(\mathbf{z},0)^*A(\mathbf{z},0)=\mathbf{z}^*B\mathbf{z}$.\\
For the lower bound, set $\mathcal{S}=\operatorname{span}((\mathbf{v}_1,0),\ldots,(\mathbf{v}_i,0))$ and $\mathcal{T}=\operatorname{span}(\mathbf{u}_i,\ldots,\mathbf{u}_{n+1})$. Their dimensions are $i$ and $n+2-i$, so $\dim\mathcal{S}+\dim\mathcal{T}>n+1$ and $\mathcal{S}\cap\mathcal{T}$ contains a nonzero vector.\\
Choose a unit $\mathbf{x}=(\mathbf{z},0)$ in the intersection. The spectral expansion in the $\mathbf{u}_j$ gives $\mathbf{x}^*A\mathbf{x}\geq\lambda_i(A)$, while the expansion in the $\mathbf{v}_j$ gives $\mathbf{z}^*B\mathbf{z}\leq\lambda_i(B)$. Hence $\lambda_i(A)\leq\lambda_i(B)$.\\
For the upper bound, take $\mathcal{S}'=\operatorname{span}((\mathbf{v}_i,0),\ldots,(\mathbf{v}_n,0))$ and $\mathcal{T}'=\operatorname{span}(\mathbf{u}_1,\ldots,\mathbf{u}_{i+1})$. Again the dimensions sum to $n+2$, so choose a unit $\mathbf{x}=(\mathbf{z},0)\in\mathcal{S}'\cap\mathcal{T}'$.\\
Now $\mathbf{z}^*B\mathbf{z}\geq\lambda_i(B)$ and $\mathbf{x}^*A\mathbf{x}\leq\lambda_{i+1}(A)$, giving $\lambda_i(B)\leq\lambda_{i+1}(A)$.\\
\hlt{Structural Commentary}: A principal submatrix is a coordinate slice of the larger quadratic form.\\
\hlt{Key Insight}: Deleting one coordinate costs one dimension, so eigenvalues shift by at most one index.
\end{proof}
\end{theorem}

\begin{corollary} \label{cor:singular-values-interlace-after-deleting-row-or-column} \hlt{Singular Values Interlace After Deleting a Row or Column}\\
Let $B$ be obtained from $A\in\operatorname{Mat}(n,n,\mathbb{C})$ by deleting one row or column.\\
With singular values ordered decreasingly, $\sigma_i(A)\geq\sigma_i(B)\geq\sigma_{i+1}(A)\quad(1\leq i<n)$.
\begin{proof}
For a row deletion, permute rows so that $PA=\begin{bmatrix}B\\ \mathbf{y}^*\end{bmatrix}$. \\
Then $PA$ has the same singular values as $A$, and $BB^*$ is a principal submatrix of $(PA)(PA)^*$.\\
Cauchy interlacing applied to these Gram matrices gives the result after taking square roots and reversing the increasing eigenvalue order. A column deletion follows by applying the row case to $A^*$.\\
\hlt{Structural Commentary}: Row deletion makes $BB^*$ a principal submatrix of a permuted $AA^*$.\\
\hlt{Key Insight}: Singular-value interlacing is Cauchy interlacing applied to Gram matrices.
\end{proof}
\end{corollary}

\begin{corollary} \label{cor:cauchy-interlacing-general-principal-submatrix} \hlt{Cauchy Interlacing for Principal Submatrices}\\
Let $B\in\operatorname{Mat}(n,n,\mathbb{C})$ be a principal submatrix of Hermitian $A\in\operatorname{Mat}(n+m,n+m,\mathbb{C})$.\\
Then, for $1\leq i\leq n$, $\lambda_i(A)\leq\lambda_i(B)\leq\lambda_{i+m}(A)$.
\begin{proof}
After a permutation similarity, $B$ is leading.\\
Delete the other $m$ coordinates one at a time and apply Theorem \ref{thm:cauchy-principal-submatrix-interlacing} at each step.\\
\hlt{Structural Commentary}: Principal submatrices come from repeated coordinate deletion.\\
\hlt{Key Insight}: Each deletion shifts eigenvalue indices by at most one.
\end{proof}
\end{corollary}

\begin{figure}[H]
  \centering
  \begin{tikzpicture}[
  >=Stealth,
  scale=1,
  axis/.style={->, gray!65, line width=0.65pt},
  adot/.style={circle, fill=blue!65!black, inner sep=1.8pt},
  bdot/.style={circle, fill=red!70!black, inner sep=1.8pt},
  cdot/.style={circle, fill=green!45!black, inner sep=1.8pt},
  guide/.style={dashed, gray!55, line width=0.55pt},
  lab/.style={font=\small, align=center, fill=white, inner sep=1pt},
  tiny/.style={font=\small, align=center, fill=white, inner sep=1pt}
]
  \draw[axis] (0,1.05) -- (3.55,1.05) node[right, lab] {$A+E$};
  \draw[axis] (0,0) -- (3.55,0) node[right, lab] {$A$};

  \foreach \x in {0.85,2.55} {
    \node[adot] at (\x,0) {};
    \draw[guide] (\x,1.22) -- (\x,-0.24);
  }
  \node[bdot] at (1.45,1.05) {};
  \draw[<->, blue!55!black] (0.85,0.34) -- (2.55,0.34);
  \node[tiny, below=5pt] at (0.85,0) {$\lambda_i(A)$};
  \node[tiny, below=5pt] at (2.55,0) {$\lambda_{i+1}(A)$};
  \node[tiny, above=5pt] at (1.45,1.05) {$\lambda_i(A+E)$};

  \draw[axis] (5.1,1.05) -- (8.95,1.05) node[right, lab] {$B$};
  \draw[axis] (5.1,0) -- (8.95,0) node[right, lab] {$A$};

  \foreach \x in {6.0,7.9} {
    \node[adot] at (\x,0) {};
    \draw[guide] (\x,1.22) -- (\x,-0.24);
  }
  \node[cdot] at (6.65,1.05) {};
  \draw[<->, green!45!black] (6.0,0.34) -- (7.9,0.34);
  \node[tiny, below=5pt] at (6.0,0) {$\lambda_i(A)$};
  \node[tiny, below=5pt] at (7.9,0) {$\lambda_{i+m}(A)$};
  \node[tiny, above=5pt] at (6.65,1.05) {$\lambda_i(B)$};
\end{tikzpicture}

  \caption{Rank-one and principal-submatrix interlacing place new eigenvalues between old ones.}
  \label{fig:eigenvalue-interlacing-ladder}
\end{figure}

\begin{theorem} \label{thm:schur-majorization-eigenvalues-by-diagonal} \hlt{Schur Majorization of the Diagonal}\\
Let $A\in\operatorname{Mat}(n,n,\mathbb{C})$ be Hermitian, with increasing eigenvalues $\lambda_i(A)$ and increasing diagonal entries $d_i$.\\
Then $\sum\limits_{i=1}^{k}\lambda_i(A)\leq \sum\limits_{i=1}^{k}d_i$, $(1\leq k\leq n)$ with equality for $k=n$.
\begin{proof}
Permute coordinates so that $a_{ii}=d_i$, and let $A_k$ be the leading $k\times k$ principal submatrix.\\
Cauchy interlacing gives $\lambda_i(A)\leq\lambda_i(A_k)$ for $1\leq i\leq k$, hence $\sum\limits_{i=1}^{k}\lambda_i(A)\leq\sum\limits_{i=1}^{k}\lambda_i(A_k)=\operatorname{tr}(A_k)=\sum\limits_{i=1}^{k}d_i$.\\
For $k=n$, this is $\operatorname{tr}(A)=\operatorname{tr}(A)$.\\
\hlt{Structural Commentary}: Ordered diagonal sums are traces of leading principal blocks.\\
\hlt{Key Insight}: Diagonal majorization is Cauchy interlacing plus trace.
\end{proof}
\end{theorem}

\begin{method} \label{meth:eigenvalue-interlacing-method} \hlt{Eigenvalue Interlacing Method}\\
To prove or use an interlacing estimate:
\begin{enumerate}[label=(\roman*)]
\setlength{\itemsep}{0pt}
\item order Hermitian eigenvalues increasingly;
\item identify dimension lost: rank-one perturbation, one deleted coordinate, $m$-coordinate principal submatrix;
\item choose test vectors by intersecting spectral subspaces with the unaffected subspace;
\item convert the Rayleigh quotient bounds into eigenvalue inequalities.
\end{enumerate}
\end{method}
